%
% Documento: Preâmbulo
%
\usepackage{lmodern}
\usepackage [latin1]{inputenc}
\renewcommand{\familydefault}{\rmdefault}
\renewcommand{\lstlistingname}{Código}

\counterwithin{quadro}{chapter}
\counterwithin{figure}{chapter}
\counterwithin{table}{chapter}
% Altera o rótulo a ser usando no elemento pré-textual "Lista de código"
\renewcommand{\lstlistlistingname}{Lista de códigos}

% Configura a ``Lista de Códigos'' conforme as regras da ABNT (para abnTeX2)
\begingroup\makeatletter
\let\newcounter\@gobble\let\setcounter\@gobbletwo
\globaldefs\@ne \let\c@loldepth\@ne
\newlistof{listings}{lol}{\lstlistlistingname}
\newlistentry{lstlisting}{lol}{0}
\endgroup

\renewcommand{\cftlstlistingaftersnum}{\hfill--\hfill}

\let\oldlstlistoflistings\lstlistoflistings
\renewcommand{\lstlistoflistings}{%
	\begingroup%
	\let\oldnumberline\numberline%
	\renewcommand{\numberline}{\lstlistingname\space\oldnumberline}%
	\oldlstlistoflistings%
	\endgroup}

\lstdefinestyle{master-thesis}{
basicstyle=\ttfamily,
numbers=left,
numberstyle=\tiny,
frame=tb,
columns=fullflexible,
showstringspaces=false
}
\lstset{style=master-thesis}

\instituicao{Universidade Federal de Minas Gerais}
\localapresentacao{Belo Horizonte - Minas Gerais}
\programa{Graduação em Engenharia de Sistemas}
\modalidade{}
\nomeautor{André Versiani de Mattos}
\titulotb{Caracterização e Extração de Atributos de Faltas de Alta Impedância em Vegetação: Aplicação em Dados Reais}
%\subtitulo{Subtítulo do trabalho}
\anoapresentacao{2026}
\grau{Bacharel em Engenharia de Sistemas}
\dataapresentacao{2026}
\mesapresentacao{Dezembro}

%Dados Orientador
\orientador{Gabriela Nunes}
\instOrientador{Cesar School}
\titulacaoorientador{Prof.ª Dra.}

%Dados Coorientador
% \coorientador{Professor Fulano}
% \instCoorientador{Instituição de Ensino}
% \titulacaocoorientador{Prof. Dr.}